\documentclass[10pt,a4paper]{amsart}
\usepackage[margin=19mm]{geometry}
\usepackage{amsmath,amssymb,mathtools}
\usepackage[scr=rsfs,cal=euler]{mathalfa}
\usepackage{xfrac}
\usepackage[unicode,hidelinks,bookmarks=false]{hyperref}
\newcommand{\ie}{i.e.\ }
\newcommand{\cf}{cf.\ }
\newcommand{\resp}{respectively\ }
\newcommand{\Subsection}{\subsection}
\providecommand{\nobreakdash}{\nobreak\hbox{-}}
\setlength{\emergencystretch}{2em}
\multlinegap0pt
\usepackage{amssymb}
\usepackage{bm}
\usepackage{tikz}
\newtheorem{lem}{Lemma}
\def\e#1\e{\begin{equation}#1\end{equation}}
\def\eq{\eqref}
\def\es{\emptyset}
\def\l{\label}
\def\le{\leqslant\nobreak}
\def\ps{\psi}
\def\sb{\subseteq}
\def\si{\sigma}
\title{Nearby Special Lagrangians}
\author{Mohammed Abouzaid, Yohsuke Imagi}
\date{Source: arXiv:2112.10385v3 (CC BY 4.0)}
\begin{document}
\maketitle

\noindent\emph{Source and licence.} Nearby Special Lagrangians. Version arXiv:2112.10385v3 (CC BY 4.0), distributed by arXiv under the Creative Commons Attribution 4.0 International licence, http://creativecommons.org/licenses/by/4.0/, as recorded in the arXiv licence metadata for this exact version and shown on the abstract page https://arxiv.org/abs/2112.10385v3. Full licence notice: \texttt{licenses/arxiv-2112.10385-CC-BY-4.0.txt}.\par\smallskip

\noindent\textit{Editorial scope.} Counted unit: the article's lem lem (source0.tex lines 2280--2284). The article's complete original proof of it is delivered with it (source0.tex lines 2285--2331, one \texttt{proof} environment of 47 lines). No mathematics is written, repaired or re-derived: the statement and the proof are the article's own lines, in the article's own notation, and no retained line is substituted or re-flowed.  No further source-local context is needed: every cross-reference of every retained line resolves inside the two delivered spans.  Nothing was omitted from any retained span, so the package carries no omission record.  No retained line contains a citation, so the wrapper carries no bibliography and no bibliography entry.  No retained line contains a figure environment, an image-inclusion command or drawing code, so no image is delivered and the package declares no asset dependency.  Editorial formatting only, in addition: an AMS \texttt{amsart} preamble that restates the article's own declarations and macros used by the retained lines, namely the environments \texttt{lem} on separate counters, together with the standard packages those lines need.  The retained environments print in this document as Lemma 1 in order of appearance, whereas the article prints the same items with the numbering of its own full document.  The complete native source of the article as distributed in the arXiv e-print for this exact version is bundled unchanged as \texttt{sources/119891/source0.tex}.
\begin{lem}\l{lem: Sob}
There exists $M_1>0$ independent of the finite cover $f:P\to Q$ and so large that for every $u\in L^2_{k+1}(T^*P)$ we have
\e\l{Sob}
\sup_P|u|_{C^1}\le M_1\|u\|_{L^2_{k+1}}\le  (M_1)^2(\|u\|_{L^2}+\|(d+ d^*f^*\ps)u\|_{L^2_k}).\e
\end{lem}

\begin{proof}
Let $\rho>0$ be less than half the injectivity radius of the Riemannian manifold $(Q,g).$ For each point $a\in Q$ denote by $B_a\subset Q$ the open geodesic ball of radius $\rho$ about $a.$
As $Q$ is compact there exists a finite set $A\subset Q$ such that $\bigcup_{a\in A}B_a=Q.$
By the Sobolev embedding on each $B_a,$ $a\in A,$ there exists $M_a>0$ independent of $f:P\to Q$ and so large that for every $u\in L^2_{k+1}(T^*B_a)$ we have
$\sup_{B_a}|u|_{C^1}\le M_a\|u\|_{L^2_{k+1}(B_a)}.$
Note that $M:=\max_{a\in A}M_a$ is also independent of $f:P\to Q.$ 
We look now at the finite cover $f:P\to Q.$ Take any $x\in P$ with $f(x)\in A$ and denote by $B_x\subset P$ the lift of $B_{f(x)}\subset Q$ with respect to $f:P\to Q,$ which exists because $B_{f(x)}$ is simply connected.
Then for any $u\in L^2_{k+1}(T^*P)$ we have
\e\l{loc Sob}|u|_{B_x}|_{C^1(B_x)}=|f_*(u|_{B_x})|_{C^1(B_a)}\le M_a\|f_*(u|_{B_x})\|_{L^2_{k+1}(B_a)}= M_a\|u|_{B_x}\|_{L^2_{k+1}(B_x)}\e
which is bounded above by $M\|u\|_{L^2_{k+1}(P)}.$ Since $P$ is covered by $B_x$'s with $f(x)\in A$ it follows by \eq{loc Sob} that $\|u\|_{L^2_{k+1}(P)}\le M\|u\|_{L^2_{k+1}(P)},$ proving the left inequality of \eq{Sob}. 

We turn now to the other part. We show first that for $a,b\in A$ with $B_a\cap B_b\ne\es$ the non-empty intersection is necessarily connected. If there were two components, there would be two distinct geodesics connecting these two.
But $B_a,B_b$ have been chosen to have the common radius $\rho$ which is less than half the injectivity radius. So $B_a\cup B_b$ lies in a single geodesic ball, which contradicts that the two geodesic exist. Thus $B_a\cap B_b$ is connected. The union $B_a\cup B_b$ is accordingly simply connected.

As $Q$ is covered by $B_a$'s of radius $\rho>0$ there exists $\si\in(0,\rho)$ so close to $\rho$ that if we denote by $B_a'\sb B_a$ the concentric ball of radius $\si$ then $(B_a')_{a\in A}$ still covers $Q.$ 
Then for $u\in L^2_{k+1}(T^*P)$ we have
\e\l{sum}
\|u\|_{L^2_{k+1}}
=\sum_{y\in f^*A} \|f_*(u|_{B_y})\|_{L^2_{k+1}(B_{f(y)})}
\le \sum_{y\in f^*A}\sum_{B_a'\cap B_{f(y)}\ne\es}\|f_*(u|_{B_y})\|_{L^2_{k+1}(B_{f(y)}\cap B_a')}.\e
For $a\in A$ and $x\in f^{-1}(a)$ denote by $Y_{ax}\sb f^*A$ the set of $y$ such that $B_{f(y)}\cap B_a'\ne\es$ and the simply connected set $B_{f(y)}\cap B_a$ lifts to $B_y\cap B_x.$
Exchanging the order of the two sums in \eq{sum} we see then that
\e\l{sumsumsum}
\begin{split}
\|u\|_{L^2_{k+1}}
&=\sum_{a\in A}\sum_{x\in f^{-1}(a)}\sum_{y\in Y_{ax}}\|f_*(u|_{B_y})\|_{L^2_{k+1}(B_{f(y)}\cap B_a')}\\
&=\sum_{a\in A}\sum_{x\in f^{-1}(a)}\sum_{y\in Y_{ax}}\|f_*(u|_{B_y\cup B_x})\|_{L^2_{k+1}(B_{f(y)}\cap B_a')}\\
&\le \sum_{a\in A}\sum_{x\in f^{-1}(a)}(\# Y_{ax})\|f_*(u|_{B_x})\|_{L^2_{k+1}(B_a')}.
\end{split}
\e
Note that there is an injection $Y_{ax}\to\{z\in A: B_a'\cap B_z\ne\es\}$ defined by $y\mapsto f(y)$ so that the number $\#Y_{ax}$ is bounded above by some $f$-independent constant $N>0.$
It follows then from \eq{sumsumsum} that
\e\l{sumsum}
\|u\|_{L^2_{k+1}}\le N\sum_{a\in A}\sum_{x\in f^{-1}(a)}\|f_*(u|_{B_x})\|_{L^2_{k+1}(B_a')}.
\e
On the other hand, by the ordinary a priori estimates over $Q$ there exists $M>0$ independent of $f:P\to Q$ and so large that for every $a\in A$ and $x\in f^{-1}(a)$ we have
\[\|f_*(u|_{B_x'})\|_{L^2_{k+1}(B_a')}\le M(\|f_*(u|_{B_x})\|_{L^2(B_a)}+\|(d+d^*\ps)f_*(u|_{B_x})\|_{L^2_k(B_a)}).\]
Hence it follows by \eq{sumsum} that
\[\begin{split}
\|u\|_{L^2_{k+1}}
&\le MN\sum_{a\in A}\sum_{x\in f^{-1}(a)} (\|f_*(u|_{B_x})\|_{L^2(B_a)}+\|(d+d^*\ps)f_*(u|_{B_x})\|_{L^2_k(B_a)})\\
&\le MN\sum_{a\in A}\sum_{x\in f^{-1}(a)} 
(\|u\|_{L^2(B_x)}+\|(d+d^*f^*\ps)u\|_{L^2_k(B_x)})\\
&= MN(\|u\|_{L^2(P)}+\|(d+d^*f^*\ps)u\|_{L^2_k(P)}).
\end{split}\]
Thus $M_1:=MN$ will do for the right part of \eq{Sob}.
\end{proof}

\end{document}
